% Generated draft: requires mathematical and source review.
\documentclass[11pt]{article}
\usepackage[margin=25mm]{geometry}
\usepackage{fontspec}
\setmainfont{Latin Modern Roman}
\usepackage{amsmath,amssymb,mathtools,unicode-math}
\setmathfont{Latin Modern Math}
\usepackage{enumitem,needspace,xurl,hyperref}
\setlength{\parindent}{0pt}
\setlength{\parskip}{5pt}
\setlength{\emergencystretch}{3em}
\newcommand{\N}{\mathbb{N}}
\newcommand{\Z}{\mathbb{Z}}
\newcommand{\Q}{\mathbb{Q}}
\newcommand{\R}{\mathbb{R}}
\newcommand{\C}{\mathbb{C}}
\newcommand{\F}{\mathbb{F}}
\newcommand{\A}{\mathbb{A}}
\newcommand{\PP}{\mathbb{P}}
\newcommand{\E}{\mathbb{E}}
\newcommand{\eps}{\varepsilon}
\newcommand{\dd}{\,\mathrm{d}}
\begin{document}
\section*{UnsolvedMath ID: 30002301}
\textbf{Source problem code:} catalogue entry
\subsection{Definitions and mathematical statement}
The supplied record does not explicitly identify the underlying function space or the translation convention. The following formulation uses the conventional real-line interpretation, which requires confirmation against the cited report: let \(H=L^2(\R;\C)\), with inner product linear in the first variable, and define
\[
W(g)=\{g_{j,k}:j,k\in\Z\},\qquad g_{j,k}(x)=2^{j/2}g(2^j x-k).
\]
The system \(W(g)\) is Bessel if there exists \(B<\infty\) such that
\[
\sum_{j,k\in\Z}|\langle f,g_{j,k}\rangle|^2\le B\|f\|_2^2\qquad(f\in H).
\]
For \(g_1,g_2\in H\), let \(W(g_1,g_2)\) denote the indexed family \((g_{r,j,k})_{r\in\{1,2\},\,j,k\in\Z}\), where \(g_{r,j,k}(x)=2^{j/2}g_r(2^j x-k)\). It is a frame if there exist \(0<A\le B<\infty\) such that
\[
A\|f\|_2^2\le\sum_{r=1}^2\sum_{j,k\in\Z}|\langle f,g_{r,j,k}\rangle|^2\le B\|f\|_2^2\qquad(f\in H).
\]
Two such frames \(W(g_1,g_2)\) and \(W(h_1,h_2)\) are dual if, for every \(f\in H\),
\[
f=\sum_{r=1}^2\sum_{j,k\in\Z}\langle f,h_{r,j,k}\rangle g_{r,j,k}
 =\sum_{r=1}^2\sum_{j,k\in\Z}\langle f,g_{r,j,k}\rangle h_{r,j,k},
\]
with unconditional convergence in \(H\).

\textbf{Completion question under this interpretation.} For every pair \(\psi_1,\widetilde\psi_1\in H\) for which \(W(\psi_1)\) and \(W(\widetilde\psi_1)\) are Bessel, do there exist \(\psi_2,\widetilde\psi_2\in H\) such that \(W(\psi_1,\psi_2)\) and \(W(\widetilde\psi_1,\widetilde\psi_2)\) are dual frames for \(H\)? The original generators are retained, and one additional generator is sought for each system. No Fourier-support or multiresolution hypotheses are imposed in the supplied unrestricted question.

\textbf{Unresolved alternative formulation.} The record also describes this question as equivalent to Han's conjectured tight-frame completion for every sufficiently large bound \(K\). A frame is tight with bound \(K>0\) when its coefficient-square sum equals \(K\|f\|_2^2\) for every \(f\in H\). However, the record does not give the full statement attributed to Han, the precise completion involved, or the threshold and its dependence on the initial data. Consequently, the claimed equivalence and its exact quantifiers are retained as supplied qualifications, not asserted here as established mathematical facts.
\subsection{Short English statement}
Under an explicitly provisional \(L^2(\R;\C)\) interpretation, ask whether every pair of Bessel dyadic wavelet systems can each receive one additional generator so that the enlarged systems are dual frames. The record also attributes an equivalent tight-frame completion question to Han, but does not specify it fully. The supplied open status is dated curation, not an independently verified conclusion.
\subsection{Sources}
\begin{itemize}
\item Supplied local UnsolvedMath registry record, permanent ID \(30002301\), source problem code \texttt{catalogue entry}, title \emph{Dual-Frame Completion of Bessel Wavelet Systems}; collection: Oberwolfach Reports Open Problems; proposed year: 2013. No proposer is identified. The alternative conjecture is attributed to Han without a full name or a separate citation.
\item \emph{Structured Function Systems and Applications} (2013), Oberwolfach Reports, pp.~579--655, DOI: \url{https://doi.org/10.4171/owr/2013/11}. This is the supplied workshop-report provenance. The registry labels the statement \texttt{corrected\_verified} and reports comparison with the official report; that verification is supplied metadata and has not been independently repeated here.
\item \emph{Bessel sequences, wavelet frames, duals and extensions}: \url{https://doi.org/10.1016/j.indag.2018.01.009}.
\item \emph{On Parseval wavelet frames with two or three generators}: \url{https://doi.org/10.4153/CMB-2013-015-9}.
\item Literature assessment checked on \texttt{2026-08-21}: the supplied curation labels the unrestricted extension problem open and describes later results as restricted extensions, especially under Fourier-support or multiresolution hypotheses. This assessment is supplied dated curation derived from the Oberwolfach material, not independently verified here; its time-sensitive status and cited claims require independent checking.
\end{itemize}
\end{document}
